%% Un elemento K: u_h vive en los nodos (lo que ve la malla), ũ dentro del elemento (lo que no).
%% Inspirado en ~/imagenes/elemento_subescala.pdf, rehecho con la paleta de la charla y fondo transparente.
\begin{tikzpicture}[scale=0.95,line join=round]
  \coordinate (A) at (0,0); \coordinate (B) at (2.6,0.75); \coordinate (C) at (0.8,2.5);
  \coordinate (N1) at (-1.3,1.6); \coordinate (N2) at (1.9,-1.2); \coordinate (N3) at (3.0,2.6);
  \coordinate (N4) at (-1.0,-1.0); \coordinate (N5) at (4.0,0.2); \coordinate (N6) at (-0.3,3.4);
  % vecinos
  \draw[tenue!60,line width=0.4pt] (A) -- (N1) -- (C) (A) -- (N4) -- (N1) (A) -- (N2) -- (B) (N4) -- (N2)
    (B) -- (N5) -- (N3) (B) -- (N3) -- (C) (C) -- (N6) -- (N1) (N2) -- (N5) (N6) -- (N3);
  % el elemento K
  \fill[ingSub!14] (A) -- (B) -- (C) -- cycle;
  \draw[tinta,line width=1pt] (A) -- (B) -- (C) -- cycle;
  \node[font=\small] at (1.35,0.95) {$K$};
  % ũ dentro del elemento
  \foreach \p/\a in {(1.0,1.45)/20,(1.6,1.2)/140,(1.15,0.75)/260} {
    \draw[ingSub!85!black,line width=0.8pt,-{Stealth[length=3.5pt]}] \p ++(\a:0.22) arc[start angle=\a,end angle=\a+250,radius=0.22];}
  \node[font=\small,text=ingSub!85!black] at (1.75,1.75) {$\tilde{\boldsymbol u}$};
  % u_h en los nodos
  \foreach \v/\pos in {A/below left,B/right,C/above} {\fill[ingEsq] (\v) circle (2.4pt);
    \node[\pos=1pt,font=\scriptsize,text=ingEsq] at (\v) {$\boldsymbol u_h$};}
  \foreach \v in {N1,N2,N3,N4,N5,N6} \fill[ingEsq!45] (\v) circle (1.5pt);
\end{tikzpicture}
